Klassiske numeriske simulatorer er ofte noejagtige, men de kan samtidig vaere beregningstunge. Scientific Machine Learning soeger at kombinere den fysiske struktur fra differentialligninger med fleksibiliteten fra moderne maskinlaering. I dette projekt er det relevante spoergsmaal derfor ikke, om machine learning skal erstatte fysikken, men hvordan laeringsmodeller kan udnytte baade data og fysisk viden til at approksimere plasmaets dynamik.

Den mest direkte datadrevne tilgang er supervision paa simulerede data. Her traenes en model til at approksimere en funktion
\[
\hat{\mathbf{u}}_\theta(x,y,t) \approx \mathbf{u}(x,y,t),
\]
hvor maalet er at minimere forskellen mellem modelens prediktion og de felter, som simulatoren producerer. En ren datamodel kan i princippet give gode interpolationer paa kendte dynamikker, men uden yderligere struktur er der ingen garanti for, at modellen respekterer de underliggende fysiske love.

Physics-Informed Neural Networks, forkortet PINNs, adresserer dette problem ved at indbygge differentialligningerne direkte i traeningen. En PINN tager typisk rum-tidskoordinater som input og returnerer de relevante felter som output. Ved hjaelp af automatisk differentiation kan de noedvendige afledte beregnes direkte gennem netvaerket, og PDE-residualen kan derefter indgaa i tabsfunktionen. En typisk samlet tabsfunktion kan skrives som
\[
\mathcal{L} = \lambda_{\text{data}}\mathcal{L}_{\text{data}}
+ \lambda_{\text{PDE}}\mathcal{L}_{\text{PDE}}
+ \lambda_{\text{BC}}\mathcal{L}_{\text{BC}}
+ \lambda_{\text{IC}}\mathcal{L}_{\text{IC}},
\]
hvor de enkelte led maaler henholdsvis dataafvigelse, brud paa PDE-systemet, brud paa randbetingelser og brud paa begyndelsesbetingelser. Ideen er, at modellen ikke kun skal passe observerede data, men ogsaa producere loesninger, som er fysisk konsistente.

En begraensning ved klassiske PINNs er, at de ofte laerer en enkelt loesning ad gangen og kan vaere dyre at traene paa store eller komplekse systemer. Derfor er neural operators blevet et vigtigt alternativ. I stedet for at laere en funktion fra koordinater til feltvaerdier laerer en neural operator en afbildning mellem funktioner. Formelt kan man opfatte maalet som at laere en operator
\[
\mathcal{G}: a \mapsto u,
\]
hvor inputtet $a$ eksempelvis kan beskrive en begyndelsestilstand, en parameterkonfiguration eller et tidligere tidsstep, og outputtet er hele eller dele af den resulterende loesning. Denne tilgang er attraktiv, fordi den potentielt kan generalisere mellem forskellige initialbetingelser og parameterregimer.

En konkret og udbredt neural operator er Fourier Neural Operator, FNO, som udnytter konvolutionelle operationer i Fourier-rummet til effektivt at modellere globale rumlige koblinger. For PDE-problemer er dette saerligt relevant, fordi mange fysiske processer indeholder ikke-lokale sammenhaenge og multiskalaadfaerd. Naar neural operators kombineres med fysikbaserede tabsled, taler man ofte om Physics-Informed Neural Operators, PINO. Her forenes operatorlaeringens generaliseringsevne med den fysiske regularisering kendt fra PINNs.

For dette projekt er disse metodefamilier relevante, fordi de matcher problemets struktur. Der findes baade numerisk genererede reference-data fra BOUT-HESEL og et kendt PDE-system, som en model boer respektere. Det giver et naturligt grundlag for hybridmodeller, hvor data bruges til at laere dynamiske sammenhaenge, mens fysikken bruges til at afgraense, hvilke loesninger der er plausible. Det konkrete valg af modelarkitektur, inputs, outputs og traeningsprocedure beskrives senere i metodekapitlet.
